\documentclass[11pt,a4paper]{article}
\usepackage[utf8]{inputenc}
\usepackage[T1]{fontenc}
\usepackage[catalan]{babel}
\usepackage{lmodern}
\usepackage[sfdefault,lining]{FiraSans}
\usepackage[a4paper,top=2.3cm,bottom=2.2cm,left=1.8cm,right=1.8cm,headheight=1.3cm,headsep=0.4cm,footskip=1cm]{geometry}
\usepackage{amsmath,amssymb}
\usepackage{icomma}
\usepackage{graphicx}
\usepackage[table]{xcolor}
\usepackage{tikz}
\usepackage{fontawesome5}
\usepackage{ccicons}
\usepackage{fancyhdr}
\usepackage{lastpage}
\usepackage{tasks}
\usepackage{needspace}
\usepackage[most]{tcolorbox}
\usepackage{array}
\usepackage[hidelinks,pdftitle={Fitxa d'exercicis de la unitat 10: Distribucions de probabilitat. 1r de Batxillerat CCSS},pdfsubject={Matemàtiques Aplicades a les CCSS I},pdfkeywords={variable aleatòria, funció de probabilitat, funció de distribució, esperança, distribució binomial, funció de densitat, distribució uniforme, distribució normal, normal típica, tipificació, aproximació de la binomial a la normal}]{hyperref}

\definecolor{navy}{HTML}{1B2748}
\definecolor{accent}{HTML}{E8603A}
\definecolor{teal}{HTML}{1F9E90}
\definecolor{soft}{HTML}{EEF1F8}
\definecolor{muted}{HTML}{5D6685}

\setlength{\parindent}{0pt}
\setlength{\parskip}{3pt}
\settasks{label=\alph*), label-format=\bfseries\color{navy}, label-width=1.3em, item-indent=1.8em, column-sep=0.6em, after-item-skip=5pt, before-skip=3pt, after-skip=4pt}

% ---------- capçalera i peu ----------
\newcommand{\llicencia}{\href{https://creativecommons.org/licenses/by-nc-sa/4.0/deed.ca}{\ccbyncsa}\ \ Llicència Creative Commons Reconeixement-NoComercial-CompartirIgual 4.0 Internacional (CC BY-NC-SA 4.0)}
\pagestyle{fancy}
\fancyhf{}
\fancyhead[L]{\raisebox{-0.35\height}{\includegraphics[height=0.95cm]{logo-gva.png}}}
\fancyhead[C]{\small\color{muted}Fitxa d'exercicis · Unitat 10: Distribucions de probabilitat · 1r Batxillerat CCSS}
\fancyhead[R]{\raisebox{-0.35\height}{\includegraphics[height=0.95cm]{logo-vermellar.png}}}
\fancyfoot[L]{\scriptsize\color{muted}\llicencia}
\fancyfoot[R]{\small\color{muted}\thepage}
\renewcommand{\headrulewidth}{0.4pt}
\renewcommand{\headrule}{\color{muted!40}\hrule height \headrulewidth}
\fancypagestyle{primera}{\fancyhf{}\renewcommand{\headrulewidth}{0pt}
  \fancyfoot[L]{\scriptsize\color{muted}\llicencia}
  \fancyfoot[R]{\small\color{muted}\thepage}}

% ---------- nivells i exercicis ----------
\newcommand{\estrelles}[1]{{\color{accent}\ifcase#1\or\faStar\faStar[regular]\faStar[regular]\or\faStar\faStar\faStar[regular]\or\faStar\faStar\faStar\fi}}
\newcounter{exer}
\newcommand{\ex}[1]{\par\needspace{4\baselineskip}\medskip\refstepcounter{exer}\noindent{\color{navy}\bfseries\theexer.}\ \ }

\newcommand{\seccio}[2]{\par\needspace{16\baselineskip}\bigskip
  \begin{tcolorbox}[colback=navy,colframe=navy,coltext=white,arc=2mm,boxrule=0pt,left=3mm,right=3mm,top=1.2mm,bottom=1.2mm]
  {\large\bfseries #1\quad #2}
  \end{tcolorbox}\vspace{-2pt}}

\usepackage{enumitem}
\newcommand{\exf}[1]{\par\medskip\refstepcounter{exer}\noindent{\color{navy}\bfseries\theexer.}\ \ }
\providecommand{\gt}{>}\providecommand{\lt}{<}
% quadre de recordatori (només a la fitxa d'exercicis)
\newcommand{\recorda}[1]{\par\smallskip\begin{tcolorbox}[colback=accent!7,colframe=accent!55,arc=2mm,boxrule=0.5pt,left=2.5mm,right=2.5mm,top=1mm,bottom=1mm]\small{\color{accent}\faLightbulb[regular]}\ \textbf{Recorda:}\ #1\end{tcolorbox}}
\newcommand{\pau}[1]{{\color{teal}\small\bfseries\faUniversity\ PAU · #1.}\ }
\newcommand{\res}[1]{#1}\newcommand{\rest}[1]{#1}

\def\fmneed{11}% que el títol d'una secció no quede sol al final del full
% =====================================================================
% Capçalera FentMates de les fitxes: un color per a cada curs.
% Els colors són els de la pàgina inicial (LEVELS a assets/js/materials.js);
% el logo ve de _eines/logo/pdf/ (python3 _eines/logo/genera.py).
%
% Ús, al final del preàmbul de la unitat:
%   % =====================================================================
% Capçalera FentMates de les fitxes: un color per a cada curs.
% Els colors són els de la pàgina inicial (LEVELS a assets/js/materials.js);
% el logo ve de _eines/logo/pdf/ (python3 _eines/logo/genera.py).
%
% Ús, al final del preàmbul de la unitat:
%   % =====================================================================
% Capçalera FentMates de les fitxes: un color per a cada curs.
% Els colors són els de la pàgina inicial (LEVELS a assets/js/materials.js);
% el logo ve de _eines/logo/pdf/ (python3 _eines/logo/genera.py).
%
% Ús, al final del preàmbul de la unitat:
%   \input{../comu/fentmates.tex}
%   \fmcurs{1eso}
% i, al començament del document:
%   \capcalera{Matemàtiques}{1r ESO}{Fitxa d'exercicis · Unitat 2}{Divisibilitat}
% Opcional abans de \input: \def\fmneed{14} (espai mínim abans de cada secció).
% =====================================================================
\makeatletter
\def\fm@logodir{../../logo/pdf}
\providecommand{\fmneed}{10}
\def\fm@cursnom{}

% clau del curs, color i color del text damunt del color
\newcommand{\fm@curs}[3]{\@namedef{fm@color@#1}{#2}\@namedef{fm@text@#1}{#3}}
\fm@curs{1eso}{E6B84A}{navy}
\fm@curs{2eso}{3B82F6}{white}
\fm@curs{3eso}{FF7A50}{white}
\fm@curs{4eso-a}{A78BFA}{white}
\fm@curs{4eso-b}{6D28D9}{white}
\fm@curs{1batx-mi}{9B1C1C}{white}
\fm@curs{1batx-ccss}{EF6B6B}{white}
\fm@curs{2batx-mii}{1E6B47}{white}
\fm@curs{2batx-ccss}{2BB3A3}{white}

\newcommand{\fmcurs}[1]{%
  \@ifundefined{fm@color@#1}{\PackageError{fentmates}{Curs desconegut: #1}{}}{}%
  \edef\fm@tmp{\noexpand\definecolor{curs}{HTML}{\@nameuse{fm@color@#1}}%
               \noexpand\colorlet{curstext}{\@nameuse{fm@text@#1}}}\fm@tmp
  \colorlet{cursfosc}{curs!72!black}%
  \def\fm@logo{\fm@logodir/vertical-#1.pdf}%
  \def\fm@marca{\fm@logodir/marca-#1.pdf}}

% ---------- capçalera de la primera pàgina ----------
% #1 matèria (petit), #2 curs, #3 tipus de document i unitat, #4 títol
\newcommand{\capcalera}[4]{%
  \gdef\fm@cursnom{#2}%
  \par\noindent
  \begin{tcolorbox}[colback=curs,colframe=curs,boxrule=0pt,arc=2.2mm,width=3.5cm,height=2.15cm,
      valign=center,left=3mm,right=2mm,top=1mm,bottom=1mm,nobeforeafter,box align=center]
    \color{curstext}{\footnotesize #1}\par\vspace{1pt}{\fontsize{19}{21}\selectfont\bfseries #2\par}%
  \end{tcolorbox}\hspace{5mm}%
  \begin{minipage}[c]{\dimexpr\linewidth-3.5cm-5mm-2.9cm\relax}
    {\color{cursfosc}\small\bfseries\MakeUppercase{#3}}\par\vspace{2pt}
    {\color{navy}\fontsize{20}{23}\selectfont\bfseries #4\par}
  \end{minipage}\hfill
  \begin{minipage}[c]{2.6cm}\raggedleft\includegraphics[width=2.6cm]{\fm@logo}\end{minipage}\par
  \vspace{2pt}}

% ---------- barres de secció del color del curs ----------
\renewcommand{\seccio}[2]{\par\needspace{\fmneed\baselineskip}\bigskip
  \begin{tcolorbox}[colback=curs,colframe=curs,coltext=curstext,arc=2mm,boxrule=0pt,left=3mm,right=3mm,top=1.2mm,bottom=1.2mm]
  {\large\bfseries #1\quad #2}
  \end{tcolorbox}\vspace{-2pt}}

% ---------- estrelles de dificultat del color del curs (a Batxillerat no se'n posen) ----------
\@ifundefined{estrelles}{}{%
  \renewcommand{\estrelles}[1]{{\color{curs}\ifcase#1\or\faStar\faStar[regular]\faStar[regular]\or\faStar\faStar\faStar[regular]\or\faStar\faStar\faStar\fi}}}

% ---------- capçalera i peu de les pàgines ----------
% primera pàgina: logos institucionals al peu
\fancypagestyle{primera}{\fancyhf{}\renewcommand{\headrulewidth}{0pt}
  \fancyfoot[L]{\raisebox{-0.32\height}{\includegraphics[height=0.78cm]{logo-gva.png}}\hspace{5mm}%
    \raisebox{-0.32\height}{\includegraphics[height=0.78cm]{logo-vermellar.png}}}
  \fancyfoot[R]{\scriptsize\color{muted}\href{https://creativecommons.org/licenses/by-nc-sa/4.0/deed.ca}{\ccbyncsa}\ \ CC BY-NC-SA 4.0 · fentmates.org\qquad\small\thepage}}
% la resta: el símbol a l'esquerra i el curs a la dreta, amb el títol al centre com abans
\fancyhead[L]{\raisebox{-0.3\height}{\includegraphics[height=0.62cm]{\fm@marca}}}
\fancyhead[R]{\small\bfseries\color{cursfosc}\fm@cursnom}
\renewcommand{\headrule}{\color{curs!45}\hrule height \headrulewidth}
\makeatother

%   \fmcurs{1eso}
% i, al començament del document:
%   \capcalera{Matemàtiques}{1r ESO}{Fitxa d'exercicis · Unitat 2}{Divisibilitat}
% Opcional abans de \input: \def\fmneed{14} (espai mínim abans de cada secció).
% =====================================================================
\makeatletter
\def\fm@logodir{../../logo/pdf}
\providecommand{\fmneed}{10}
\def\fm@cursnom{}

% clau del curs, color i color del text damunt del color
\newcommand{\fm@curs}[3]{\@namedef{fm@color@#1}{#2}\@namedef{fm@text@#1}{#3}}
\fm@curs{1eso}{E6B84A}{navy}
\fm@curs{2eso}{3B82F6}{white}
\fm@curs{3eso}{FF7A50}{white}
\fm@curs{4eso-a}{A78BFA}{white}
\fm@curs{4eso-b}{6D28D9}{white}
\fm@curs{1batx-mi}{9B1C1C}{white}
\fm@curs{1batx-ccss}{EF6B6B}{white}
\fm@curs{2batx-mii}{1E6B47}{white}
\fm@curs{2batx-ccss}{2BB3A3}{white}

\newcommand{\fmcurs}[1]{%
  \@ifundefined{fm@color@#1}{\PackageError{fentmates}{Curs desconegut: #1}{}}{}%
  \edef\fm@tmp{\noexpand\definecolor{curs}{HTML}{\@nameuse{fm@color@#1}}%
               \noexpand\colorlet{curstext}{\@nameuse{fm@text@#1}}}\fm@tmp
  \colorlet{cursfosc}{curs!72!black}%
  \def\fm@logo{\fm@logodir/vertical-#1.pdf}%
  \def\fm@marca{\fm@logodir/marca-#1.pdf}}

% ---------- capçalera de la primera pàgina ----------
% #1 matèria (petit), #2 curs, #3 tipus de document i unitat, #4 títol
\newcommand{\capcalera}[4]{%
  \gdef\fm@cursnom{#2}%
  \par\noindent
  \begin{tcolorbox}[colback=curs,colframe=curs,boxrule=0pt,arc=2.2mm,width=3.5cm,height=2.15cm,
      valign=center,left=3mm,right=2mm,top=1mm,bottom=1mm,nobeforeafter,box align=center]
    \color{curstext}{\footnotesize #1}\par\vspace{1pt}{\fontsize{19}{21}\selectfont\bfseries #2\par}%
  \end{tcolorbox}\hspace{5mm}%
  \begin{minipage}[c]{\dimexpr\linewidth-3.5cm-5mm-2.9cm\relax}
    {\color{cursfosc}\small\bfseries\MakeUppercase{#3}}\par\vspace{2pt}
    {\color{navy}\fontsize{20}{23}\selectfont\bfseries #4\par}
  \end{minipage}\hfill
  \begin{minipage}[c]{2.6cm}\raggedleft\includegraphics[width=2.6cm]{\fm@logo}\end{minipage}\par
  \vspace{2pt}}

% ---------- barres de secció del color del curs ----------
\renewcommand{\seccio}[2]{\par\needspace{\fmneed\baselineskip}\bigskip
  \begin{tcolorbox}[colback=curs,colframe=curs,coltext=curstext,arc=2mm,boxrule=0pt,left=3mm,right=3mm,top=1.2mm,bottom=1.2mm]
  {\large\bfseries #1\quad #2}
  \end{tcolorbox}\vspace{-2pt}}

% ---------- estrelles de dificultat del color del curs (a Batxillerat no se'n posen) ----------
\@ifundefined{estrelles}{}{%
  \renewcommand{\estrelles}[1]{{\color{curs}\ifcase#1\or\faStar\faStar[regular]\faStar[regular]\or\faStar\faStar\faStar[regular]\or\faStar\faStar\faStar\fi}}}

% ---------- capçalera i peu de les pàgines ----------
% primera pàgina: logos institucionals al peu
\fancypagestyle{primera}{\fancyhf{}\renewcommand{\headrulewidth}{0pt}
  \fancyfoot[L]{\raisebox{-0.32\height}{\includegraphics[height=0.78cm]{logo-gva.png}}\hspace{5mm}%
    \raisebox{-0.32\height}{\includegraphics[height=0.78cm]{logo-vermellar.png}}}
  \fancyfoot[R]{\scriptsize\color{muted}\href{https://creativecommons.org/licenses/by-nc-sa/4.0/deed.ca}{\ccbyncsa}\ \ CC BY-NC-SA 4.0 · fentmates.org\qquad\small\thepage}}
% la resta: el símbol a l'esquerra i el curs a la dreta, amb el títol al centre com abans
\fancyhead[L]{\raisebox{-0.3\height}{\includegraphics[height=0.62cm]{\fm@marca}}}
\fancyhead[R]{\small\bfseries\color{cursfosc}\fm@cursnom}
\renewcommand{\headrule}{\color{curs!45}\hrule height \headrulewidth}
\makeatother

%   \fmcurs{1eso}
% i, al començament del document:
%   \capcalera{Matemàtiques}{1r ESO}{Fitxa d'exercicis · Unitat 2}{Divisibilitat}
% Opcional abans de \input: \def\fmneed{14} (espai mínim abans de cada secció).
% =====================================================================
\makeatletter
\def\fm@logodir{../../logo/pdf}
\providecommand{\fmneed}{10}
\def\fm@cursnom{}

% clau del curs, color i color del text damunt del color
\newcommand{\fm@curs}[3]{\@namedef{fm@color@#1}{#2}\@namedef{fm@text@#1}{#3}}
\fm@curs{1eso}{E6B84A}{navy}
\fm@curs{2eso}{3B82F6}{white}
\fm@curs{3eso}{FF7A50}{white}
\fm@curs{4eso-a}{A78BFA}{white}
\fm@curs{4eso-b}{6D28D9}{white}
\fm@curs{1batx-mi}{9B1C1C}{white}
\fm@curs{1batx-ccss}{EF6B6B}{white}
\fm@curs{2batx-mii}{1E6B47}{white}
\fm@curs{2batx-ccss}{2BB3A3}{white}

\newcommand{\fmcurs}[1]{%
  \@ifundefined{fm@color@#1}{\PackageError{fentmates}{Curs desconegut: #1}{}}{}%
  \edef\fm@tmp{\noexpand\definecolor{curs}{HTML}{\@nameuse{fm@color@#1}}%
               \noexpand\colorlet{curstext}{\@nameuse{fm@text@#1}}}\fm@tmp
  \colorlet{cursfosc}{curs!72!black}%
  \def\fm@logo{\fm@logodir/vertical-#1.pdf}%
  \def\fm@marca{\fm@logodir/marca-#1.pdf}}

% ---------- capçalera de la primera pàgina ----------
% #1 matèria (petit), #2 curs, #3 tipus de document i unitat, #4 títol
\newcommand{\capcalera}[4]{%
  \gdef\fm@cursnom{#2}%
  \par\noindent
  \begin{tcolorbox}[colback=curs,colframe=curs,boxrule=0pt,arc=2.2mm,width=3.5cm,height=2.15cm,
      valign=center,left=3mm,right=2mm,top=1mm,bottom=1mm,nobeforeafter,box align=center]
    \color{curstext}{\footnotesize #1}\par\vspace{1pt}{\fontsize{19}{21}\selectfont\bfseries #2\par}%
  \end{tcolorbox}\hspace{5mm}%
  \begin{minipage}[c]{\dimexpr\linewidth-3.5cm-5mm-2.9cm\relax}
    {\color{cursfosc}\small\bfseries\MakeUppercase{#3}}\par\vspace{2pt}
    {\color{navy}\fontsize{20}{23}\selectfont\bfseries #4\par}
  \end{minipage}\hfill
  \begin{minipage}[c]{2.6cm}\raggedleft\includegraphics[width=2.6cm]{\fm@logo}\end{minipage}\par
  \vspace{2pt}}

% ---------- barres de secció del color del curs ----------
\renewcommand{\seccio}[2]{\par\needspace{\fmneed\baselineskip}\bigskip
  \begin{tcolorbox}[colback=curs,colframe=curs,coltext=curstext,arc=2mm,boxrule=0pt,left=3mm,right=3mm,top=1.2mm,bottom=1.2mm]
  {\large\bfseries #1\quad #2}
  \end{tcolorbox}\vspace{-2pt}}

% ---------- estrelles de dificultat del color del curs (a Batxillerat no se'n posen) ----------
\@ifundefined{estrelles}{}{%
  \renewcommand{\estrelles}[1]{{\color{curs}\ifcase#1\or\faStar\faStar[regular]\faStar[regular]\or\faStar\faStar\faStar[regular]\or\faStar\faStar\faStar\fi}}}

% ---------- capçalera i peu de les pàgines ----------
% primera pàgina: logos institucionals al peu
\fancypagestyle{primera}{\fancyhf{}\renewcommand{\headrulewidth}{0pt}
  \fancyfoot[L]{\raisebox{-0.32\height}{\includegraphics[height=0.78cm]{logo-gva.png}}\hspace{5mm}%
    \raisebox{-0.32\height}{\includegraphics[height=0.78cm]{logo-vermellar.png}}}
  \fancyfoot[R]{\scriptsize\color{muted}\href{https://creativecommons.org/licenses/by-nc-sa/4.0/deed.ca}{\ccbyncsa}\ \ CC BY-NC-SA 4.0 · fentmates.org\qquad\small\thepage}}
% la resta: el símbol a l'esquerra i el curs a la dreta, amb el títol al centre com abans
\fancyhead[L]{\raisebox{-0.3\height}{\includegraphics[height=0.62cm]{\fm@marca}}}
\fancyhead[R]{\small\bfseries\color{cursfosc}\fm@cursnom}
\renewcommand{\headrule}{\color{curs!45}\hrule height \headrulewidth}
\makeatother

\fmcurs{1batx-ccss}

\begin{document}
\thispagestyle{primera}

\capcalera{Matemàtiques CCSS}{1r Batx}{Fitxa d'exercicis · Unitat 10}{\raggedright Distribucions de probabilitat}
\vspace{10pt}
Nom i cognoms: \rule{8.4cm}{0.4pt}\hfill Grup: \rule{1.6cm}{0.4pt}\hfill Data: \rule{2.2cm}{0.4pt}
\vspace{8pt}

\begin{tcolorbox}[colback=soft,colframe=soft,arc=2mm,left=3mm,right=3mm,top=1.5mm,bottom=1.5mm]
\small Resol els exercicis al quadern mostrant tots els passos: identifica la variable aleatòria i la seua distribució, fes un dibuix de l'àrea que busques i arredoneix $z$ a dues decimals. Fes servir les taules de la binomial i de la normal $N(0,1)$ que tens als apunts interactius de la unitat (són les de la PAU valenciana). Els 12 exercicis marcats amb \textcolor{teal}{\faUniversity\ PAU} són de les proves d'accés a la universitat.
\end{tcolorbox}

% =====================================================
\seccio{1}{Variables aleatòries discretes}

\exf{1} Digues si les variables aleatòries següents són discretes o contínues:\begin{tasks}(2)
  \task El nombre de germans d'un alumne.
  \task El temps que tarda un autobús a arribar.
  \task El pes d'un nadó.
  \task El nombre de clients que entren en una botiga en una hora.
  \task La temperatura màxima d'un dia.
  \task El nombre de cares en 10 llançaments d'una moneda.
\end{tasks}

\ex{1} Llancem dues monedes i $X$ és el nombre de cares obtingudes.\begin{tasks}(1)
  \task Fes la taula de la funció de probabilitat i representa-la amb un diagrama de barres.
  \task Troba la funció de distribució $F(x)=P(X\leq x)$ per a cada valor.
  \task Calcula la mitjana i la desviació típica.
\end{tasks}\recorda{$\mu=\sum x_i\,p_i$; \ $\sigma^2=\sum x_i^2\,p_i-\mu^2$; \ $\sigma=\sqrt{\sigma^2}$.}

\ex{1} Una variable aleatòria $X$ pren els valors 1, 2, 3 i 4 amb probabilitats $k$, $2k$, $3k$ i $4k$.\begin{tasks}(3)
  \task Calcula $k$.
  \task Calcula $\mu$ i $\sigma$.
  \task Calcula $P(X\geq3)$ i $P(X\lt3)$.
\end{tasks}

\ex{2} Llancem dos daus i $X$ és la suma dels punts. Fes la taula de la funció de probabilitat i calcula $P(X\geq10)$, la mitjana i la desviació típica.

\ex{2} Una urna conté 4 boles blanques i 2 de negres. Traiem dues boles sense reemplaçament i $X$ és el nombre de boles negres. Troba la funció de probabilitat de $X$, la mitjana i la desviació típica.\recorda{fes un diagrama d'arbre: la segona extracció depén de la primera.}

\ex{2} En un joc es paga 1 euro per llançar un dau. Si ix un 6 es cobren 4 euros; si ix un 5, 2 euros; i en cas contrari no es cobra res. Siga $X$ el guany net del jugador. Calcula el guany esperat i la desviació típica. És un joc just?\recorda{el guany net és el que es cobra menys el que es paga. Un joc és just si el guany esperat és 0.}

\ex{3} Una variable aleatòria pren els valors 0, 1 i 2 amb probabilitats $a$, $0{,}5$ i $b$. Si la mitjana és $1{,}2$, calcula $a$, $b$ i la variància.

\ex{2} La funció de distribució d'una variable discreta $X$ és $F(0)=0{,}1$, $F(1)=0{,}4$, $F(2)=0{,}8$ i $F(3)=1$ (i $X$ només pren aquests valors). Troba la funció de probabilitat, $P(1\leq X\leq2)$ i la mitjana.


% =====================================================
\seccio{2}{La distribució binomial}

\exf{1} Digues si les variables següents segueixen una distribució binomial i, si és així, indica'n els paràmetres:\begin{tasks}(1)
  \task Nombre de sisos en 12 llançaments d'un dau.
  \task Nombre de boles roges en traure 3 boles, sense reemplaçament, d'una urna amb 4 roges i 6 blanques.
  \task Nombre de peces defectuoses en una mostra de 50, si el 2\,\% de la producció és defectuosa.
  \task Nombre de llançaments d'una moneda fins que ix la primera cara.
\end{tasks}\recorda{una binomial $B(n,p)$ compta els èxits en $n$ proves independents amb la mateixa probabilitat d'èxit $p$.}

\ex{1} Es llança un dau 5 vegades. Calcula la probabilitat d'obtindre exactament dos cincs i la d'obtindre almenys un cinc.\recorda{$P(X=k)=\binom nk\,p^k\,q^{\,n-k}$, \ $q=1-p$.}

\ex{2} Un test té 8 preguntes amb 3 opcions cadascuna, de les quals només una és correcta. Un alumne contesta totes les preguntes a l'atzar.\begin{tasks}(1)
  \task Quina distribució segueix el nombre d'encerts? Quants encerts s'esperen? Quina és la desviació típica?
  \task Calcula la probabilitat d'encertar-ne exactament 4 i la d'encertar-ne almenys una.
\end{tasks}

\ex{2} \pau{Castella i Lleó, PAU ordinària 2026}El 90\,\% dels joves d'entre 25 i 34 anys compra per Internet. Una botiga fa una enquesta a 10 clients joves.\begin{tasks}(1)
  \task Calcula la probabilitat que exactament 2 dels 10 clients no facen compres en línia.
  \task Calcula la probabilitat que menys de 2 clients no facen compres en línia.
\end{tasks}

\ex{1} Siga $X\sim B(5;\ 0{,}25)$. Calcula amb la taula de la binomial:\begin{tasks}(4)
  \task $P(X=2)$
  \task $P(X\leq1)$
  \task $P(X\geq4)$
  \task $P(X>0)$
\end{tasks}\recorda{la taula dona $P(X=x)$ per a $n\leq9$. Per a les acumulades, suma caselles de la mateixa columna o fes el contrari.}

\ex{2} El 80\,\% dels clients d'una botiga està satisfet amb l'atenció rebuda. Es trien 7 clients a l'atzar. Calcula, amb la taula, la probabilitat que:\begin{tasks}(1)
  \task n'estiguen satisfets exactament 6;
  \task n'estiguen satisfets almenys 6;
  \task n'estiguen satisfets com a molt 5.
\end{tasks}

\ex{1} En una família amb 4 fills, la probabilitat que cada fill siga xic és $0{,}5$. Calcula la probabilitat que hi haja exactament 2 xics, i que hi haja almenys una xica. Quants xics s'esperen?

\ex{2} Una màquina fabrica un 3\,\% de peces defectuoses. Les peces s'envasen en caixes de 50.\begin{tasks}(1)
  \task Quantes peces defectuoses hi ha, de mitjana, en una caixa? Amb quina desviació típica?
  \task Quina és la probabilitat que una caixa no tinga cap peça defectuosa? I que en tinga com a molt una?
\end{tasks}

\ex{3} Quantes vegades, com a mínim, cal llançar una moneda perquè la probabilitat d'obtindre almenys una cara siga més gran que $0{,}99$?


% =====================================================
\seccio{3}{Variables aleatòries contínues. Funció de densitat}

\exf{1} Siga $f(x)=\frac16$ si $2\leq x\leq8$ (i $f(x)=0$ en cas contrari).\begin{tasks}(1)
  \task Comprova que és una funció de densitat i dibuixa-la.
  \task Calcula $P(X\lt4)$, $P(3\leq X\leq7)$ i $P(X=5)$.
  \task Quina és la mitjana?
\end{tasks}\recorda{la probabilitat és l'àrea davall de la funció de densitat; l'àrea total ha de ser 1.}

\ex{2} Calcula $k$ perquè $f(x)=k\,x$ si $0\leq x\leq4$ (i 0 en cas contrari) siga una funció de densitat. Després calcula $P(X\leq2)$ i $P(X\gt3)$.

\ex{1} El retard d'un tren (en minuts) es distribueix uniformement entre 0 i 20 minuts. Calcula la probabilitat que el retard siga superior a 15 minuts, que estiga entre 5 i 12 minuts, i el retard mitjà.

\ex{2} La funció de densitat d'una variable $X$ és un triangle isòsceles amb base l'interval $[0,4]$ i el vèrtex en $x=2$.\begin{tasks}(1)
  \task Quina altura ha de tindre el triangle?
  \task Calcula $P(X\leq1)$, $P(1\leq X\leq3)$ i $P(X\geq3{,}5)$.
\end{tasks}

\ex{2} Raona si les funcions següents poden ser funcions de densitat (són 0 fora de l'interval indicat):\begin{tasks}(1)
  \task $f(x)=0{,}5$ en $[0,3]$
  \task $f(x)=x-1$ en $[0,2]$
  \task $f(x)=0{,}25$ en $[1,5]$
\end{tasks}


% =====================================================
\seccio{4}{La distribució normal. La normal típica $N(0,1)$ i la seua taula}

\exf{1} Si $Z\sim N(0,1)$, calcula amb la taula:\begin{tasks}(3)
  \task $P(Z\leq1{,}53)$
  \task $P(Z>2{,}05)$
  \task $P(Z<-0{,}82)$
  \task $P(Z\geq-1{,}2)$
  \task $P(0{,}75\leq Z\leq1{,}5)$
  \task $P(-2< Z<1)$
  \task $P(-2{,}58\leq Z\leq2{,}58)$
\end{tasks}\recorda{la taula dona $P(Z\leq z)$ per a $z\geq0$. \ $P(Z\geq z)=1-P(Z\leq z)$, \ $P(Z\leq-z)=1-P(Z\leq z)$. Dibuixa sempre la campana i l'àrea que busques.}

\ex{1} Comprova amb la taula la regla 68--95--99,7: calcula $P(-1\leq Z\leq1)$, $P(-2\leq Z\leq2)$ i $P(-3\leq Z\leq3)$. Què signifiquen aquests resultats per a una variable $X\sim N(\mu,\sigma)$?

\ex{2} Si $Z\sim N(0,1)$, troba el valor de $k$ en cada cas:\begin{tasks}(2)
  \task $P(Z\leq k)=0{,}8907$
  \task $P(Z\geq k)=0{,}0505$
  \task $P(Z\leq k)=0{,}25$
  \task $P(-k\leq Z\leq k)=0{,}9$
\end{tasks}\recorda{busca la probabilitat \emph{dins} de la taula; si és menor que 0,5, $k$ és negatiu.}

\ex{2} Si $Z\sim N(0,1)$, calcula $P(Z\leq-1{,}5\ \text{o}\ Z\geq1{,}5)$, \ $P(-0{,}44\leq Z\leq0)$ \ i \ $P(Z\geq-2{,}35)$. Fes un dibuix de cada àrea.


% =====================================================
\seccio{5}{Tipificació. Problemes amb la normal}

\exf{1} \pau{Illes Balears, PBAU extraordinària 2025}Quan un client entra en una botiga d'electrodomèstics, el temps (en minuts) que hi passa segueix una distribució $N(15,10)$.\begin{tasks}(1)
  \task Quina és la probabilitat que un client estiga dins de la botiga 25 minuts o menys?
  \task Quina és la probabilitat que un client estiga dins de la botiga entre 5 i 25 minuts?
\end{tasks}\recorda{si $X\sim N(\mu,\sigma)$, aleshores $Z=\dfrac{X-\mu}{\sigma}\sim N(0,1)$. Arredoneix $z$ a dues decimals i fes servir la taula.}

\ex{1} \pau{Illes Balears, PBAU ordinària 2021}El pes de les persones d'un col·legi major segueix una llei normal de mitjana 70 kg i desviació típica 15 kg. Triem una persona a l'atzar; calcula la probabilitat que:\begin{tasks}(3)
  \task pese més de 80 kg;
  \task pese menys de 50 kg;
  \task pese entre 60 i 120 kg.
\end{tasks}

\ex{1} En una universitat hi ha 2000 alumnes i el seu quocient intel·lectual (QI) segueix una distribució $N(100,15)$. Quants alumnes s'espera que tinguen un QI superior a 130? I entre 85 i 115?

\ex{2} \pau{Illes Balears, PBAU extraordinària 2023}En una població, l'altura dels homes segueix una normal de mitjana 1,76 m i desviació típica 0,12 m, i la de les dones, una normal de mitjana 1,62 m i desviació típica 0,11 m.\begin{tasks}(1)
  \task Triem un home a l'atzar. Quina és la probabilitat que la seua altura siga major o igual que 1,76 m?
  \task Triem una dona a l'atzar. Quina és la probabilitat que la seua altura siga major o igual que 1,76 m?
  \task Què és més probable, que un home mesure menys d'1,76 m o que ho faça una dona?
\end{tasks}

\ex{2} \pau{Canàries, EBAU juliol 2019}En una empresa hi ha 250 empleats. La seua edat segueix una normal de mitjana 44 anys i desviació típica 18 anys.\begin{tasks}(1)
  \task Quants empleats s'espera que tinguen més de 62 anys?
  \task Quants empleats s'espera que tinguen menys de 40 anys?
  \task Quants empleats podrien aconseguir el carnet jove de transport si el requisit és ser major d'edat i no haver complit els 30 anys?
\end{tasks}

\ex{2} \pau{Illes Balears, PAU extraordinària 2026}El temps que cada jove d'una població dedica cada dia a les xarxes socials segueix una normal de mitjana 116 minuts i desviació típica 42 minuts.\begin{tasks}(1)
  \task Quina proporció de joves dedica entre 100 i 200 minuts a les xarxes socials?
  \task En Joan és del 10\,\% de joves que més temps hi dedica. Quants minuts hi dedica, com a mínim?
\end{tasks}\recorda{en un problema invers busquem la probabilitat dins de la taula, llegim $z$ i desfem la tipificació: $x=\mu+z\,\sigma$.}

\ex{2} \pau{País Basc, EAU ordinària 2020}L'altura de les dones d'un país segueix una normal de mitjana 163 cm i desviació típica 7 cm. Una empresa de disfresses vol fer quatre talles, cadascuna adequada per al 25\,\% de les dones. Quines altures marquen el pas d'una talla a una altra?

\ex{2} \pau{País Basc, EAU ordinària 2020}El pes de les truites d'una piscifactoria segueix una normal de mitjana 250 g i desviació típica 50 g. Només són aptes per a la venda les que superen un pes determinat.\begin{tasks}(1)
  \task Quin ha de ser aquest pes si es vol que el 40\,\% de les truites siguen aptes per a la venda?
  \task Si aquest pes s'estableix en 280 g i a la piscifactoria hi ha 6000 truites, quantes se'n podran vendre?
\end{tasks}

\ex{3} \pau{Castella i Lleó, EBAU juliol 2018}En una empresa, el salari mitjà mensual és de 1400 euros i el 80,23\,\% dels treballadors cobra menys de 1570 euros. Calcula la desviació típica del salari, suposant que segueix una distribució normal.

\ex{3} \pau{La Rioja, EBAU ordinària 2023}L'estatura dels policies de França segueix una normal de desviació típica 6,5 cm, de manera que l'11,507\,\% mesura més de 183 cm.\begin{tasks}(1)
  \task Calcula la mitjana de l'estatura.
  \task Troba l'estatura que és superada pel 88,493\,\% dels policies.
\end{tasks}


% =====================================================
\seccio{6}{Aproximació de la distribució binomial a la normal}

\exf{1} Siga $X\sim B(100;\ 0{,}2)$.\begin{tasks}(1)
  \task Comprova que es pot aproximar per una normal i digues quina.
  \task Calcula aproximadament $P(X\leq15)$, $P(X=20)$ i $P(X\geq25)$.
\end{tasks}\recorda{si $np\geq5$ i $nq\geq5$, $B(n,p)\approx N\left(np,\sqrt{npq}\right)$. No oblides la correcció de continuïtat ($\pm0{,}5$).}

\ex{1} Llancem una moneda equilibrada 400 vegades. Calcula la probabilitat d'obtindre menys de 190 cares.

\ex{2} Un examen té 100 preguntes de vertader o fals. Per aprovar cal encertar-ne almenys 60. Quina és la probabilitat d'aprovar contestant totes les preguntes a l'atzar?

\ex{2} \pau{Andalusia, PAU ordinària 2026 (adaptat)}Un experiment té èxit amb probabilitat 0,4 i es repeteix 1350 vegades de manera independent. Siga $X$ el nombre d'èxits.\begin{tasks}(1)
  \task Quina distribució segueix $X$?
  \task Calcula la probabilitat que $X$ siga major que 499 i menor o igual que 580.
\end{tasks}

\ex{2} \pau{Canàries, EBAU ordinària 2022 (adaptat)}Una companyia d'assegurances té assegurats 2500 cotxes. La probabilitat que un cotxe tinga un accident al llarg d'un any és 0,1. Quants accidents s'esperen? Quina és la probabilitat que l'any que ve tinguen un accident almenys 270 cotxes?

\ex{3} Siga $X\sim B(50;\ 0{,}3)$. Calcula $P(10\leq X\leq20)$ amb l'aproximació normal. Amb un ordinador s'obté el valor exacte, $0{,}9120$. Quin error es comet?

\vfill
\begin{center}\color{muted}\small\faCheckCircle[regular]\ \ Ànim! Comprova els teus resultats amb la calculadora de probabilitats dels apunts interactius.\end{center}

\end{document}
